\RequirePackage{fix-cm}
\documentclass[11pt]{article}

\usepackage[margin=1in]{geometry}
\usepackage{amsmath,amssymb,bm}
\usepackage{array}
\usepackage{authblk}
\usepackage{booktabs}
\usepackage{graphicx}
\usepackage{pgf}
\usepackage{lmodern}
\usepackage{microtype}
\usepackage{placeins}
\usepackage[numbers,sort&compress]{natbib}
\usepackage[colorlinks=true,allcolors=blue]{hyperref}
\usepackage[nameinlink,noabbrev]{cleveref}
% Permit builds from the repository root and from paper/ (LaTeX Workshop).
\newcommand{\paperinput}[1]{\IfFileExists{paper/#1}{\input{paper/#1}}{\input{#1}}}
\newcommand{\figureinput}[1]{\IfFileExists{figures/#1}{\input{figures/#1}}{\input{../figures/#1}}}
\paperinput{defs.tex}
\graphicspath{{figures/}{../figures/}}

\title{Pressure coupling and the Biot coefficient in finite-strain poroelasticity and poroplasticity}
\author[1,2,3]{John T. Foster\thanks{Electronic address: \href{mailto:john.foster@utexas.edu}{john.foster@utexas.edu}}}
\affil[1]{The Hildebrand Department of Petroleum \& Geosystems Engineering, The University of Texas at Austin, Austin, TX 78712, USA}
\affil[2]{Department of Aerospace Engineering \& Engineering Mechanics, The University of Texas at Austin, Austin, TX 78712, USA}
\affil[3]{The Oden Institute for Computational Engineering \& Science, The University of Texas at Austin, Austin, TX 78712, USA}
\date{}
\hypersetup{
  pdftitle={Pressure coupling and the Biot coefficient in finite-strain poroelasticity and poroplasticity},
  pdfauthor={John T. Foster}
}

\begin{document}
\maketitle

\begin{abstract}
Pore pressure changes the mineral and pore volumes of a saturated porous
solid. Permanent pore-volume changes also affect its subsequent reversible
response. We derive this pressure coupling from a common stored energy that
reproduces prescribed logarithmic laws for the drained skeleton and mineral.
The construction gives a scalar mineral-volume equation and conditions for
local uniqueness. A pressure Legendre transformation identifies the reversible
Biot coefficient through both pore-volume sensitivity to deformation and total
stress sensitivity to pressure. Differentiating the mineral equation gives an
explicit coefficient and relates the effective stresses defined at fixed
pressure and fixed mineral density. The multiplier of a finite pressure change
relative to drained stress generally differs from this reversible coefficient.
Linearization recovers classical Biot elasticity.

For plastic deformation, we separate mineral shape change from permanent
pore-volume change and solve the mineral and plastic updates together.
Automatic differentiation retains their mutual dependence in the Newton
tangent. Independent checks verify mass conservation, constitutive identities,
and numerical derivatives. The finite-element formulation recovers the
analytical Mandel solution and resolves reversible coefficient changes under
larger compression. Homogeneous loading paths show how plastic history changes
the coefficient during unloading and how mineral compression affects further
plastic flow. A coupled compression test verifies plastic fluid storage and
shows a persistent coefficient increase relative to the elastic value at the
same pressure and total volume. For the synthetic material, plastic dilation
reverses fluid exchange from outflow to inflow. Material-specific experimental
validation remains future work.
\end{abstract}

\noindent\textbf{Keywords:} finite-deformation poromechanics; Biot coefficient;
poroelasticity; poroplasticity; effective stress; automatic differentiation

\paperinput{sections/introduction.tex}
\paperinput{sections/finite_deformation_biot.tex}
\paperinput{sections/implicit_ad_implementation.tex}
\paperinput{sections/poroplastic_results.tex}
\paperinput{sections/mandel_verification.tex}
\paperinput{sections/conclusions.tex}

\paperinput{sections/acknowledgements.tex}
\paperinput{sections/declarations.tex}
\paperinput{../provenance/ai_use_statement.tex}

\appendix
\paperinput{sections/mandel_analytical_appendix.tex}
\paperinput{sections/gajo_equivalence_appendix.tex}

\bibliographystyle{plainnat}
\IfFileExists{all.bib}{\bibliography{all}}{\bibliography{../all}}
\end{document}
